\documentclass[addpoints]{exam}
\usepackage{amsmath}
\usepackage{amssymb}
\usepackage{color}
\usepackage{siunitx}

% \printanswers

\newcommand{\event}{Analytic Geometry}
\newcommand{\tournament}{}
\def\type{0}

\setlength\fillinlinelength{\textwidth}
\CorrectChoiceEmphasis{\color{red}\bfseries}
\SolutionEmphasis{\color{red}\bfseries}

\setlength\answerclearance{2pt}
\pagestyle{headandfoot}
\runningheadrule
\runningheader{\event}{\tournament}{\ifnum\type=1 Class Set Test \else Full Name:\hspace{1cm} \fi}
\runningfootrule
\runningfooter{}{Page \thepage\ of \numpages}{}

\begin{document}
\begin{center}
    {\huge Grade 10 Principles of Mathematics \\ \event
    \ifprintanswers \space Key
    \else \space \ifnum\type<2 Test \fi \ifnum\type=1 \textbf{(CLASS SET)} \fi
          \ifnum\type=2 Answer Sheet \fi
    \fi \\ \vspace{3mm}\tournament\vspace{1mm}}
\end{center}

\vspace{.25\textheight}

\begin{center}
    First Name: \ifprintanswers
    \underline{\hspace{3.5cm}}\textbf{KEY}\underline{\hspace{5.5cm}}\\\ \vspace{5mm}
    \else \underline{\hspace{10cm}}\\ \vspace{5mm} \fi
    Last Name: \underline{\hspace{10cm}}\\ \vspace{5mm}
\end{center}

\noindent\textbf{\underline{Directions:}}
\begin{itemize}
    \ifnum\type=1 \item \textbf{This test is a class set.} Do not write on this paper. \fi
    \item Show all work for full marks. Leave answers in exact form where appropriate.
    \item The test is designed to be completed in 75 minutes.
\end{itemize}
\begin{center}
\textit{For grading use only}\\
\multirowgradetable{1}[pages]
\end{center}

\newpage
\begin{questions}

\section*{Part A --- Multiple Choice (15 marks)}
\vspace{5mm}

\question[1] What is the distance between the points $A(1, 2)$ and $B(4, 6)$?
\begin{choices}
    \choice $3$
    \choice $4$
    \correctchoice $5$
    \choice $\sqrt{7}$
\end{choices}
\vspace{3mm}

\question[1] What is the midpoint of the segment joining $P(-2, 4)$ and $Q(6, -2)$?
\begin{choices}
    \choice $(4, 2)$
    \correctchoice $(2, 1)$
    \choice $(1, 2)$
    \choice $(-4, 6)$
\end{choices}
\vspace{3mm}

\question[1] What is the slope of the line passing through $(3, -1)$ and $(-1, 7)$?
\begin{choices}
    \choice $2$
    \correctchoice $-2$
    \choice $\tfrac{1}{2}$
    \choice $-\tfrac{1}{2}$
\end{choices}
\vspace{3mm}

\question[1] Which equation represents a line with slope $3$ and $y$-intercept $-4$?
\begin{choices}
    \choice $3x + y = 4$
    \correctchoice $y = 3x - 4$
    \choice $y = -4x + 3$
    \choice $3y = x - 4$
\end{choices}
\vspace{3mm}

\question[1] Line $\ell_1$ has slope $\dfrac{2}{3}$. A line perpendicular to $\ell_1$ has slope:
\begin{choices}
    \choice $\dfrac{2}{3}$
    \choice $-\dfrac{2}{3}$
    \correctchoice $-\dfrac{3}{2}$
    \choice $\dfrac{3}{2}$
\end{choices}
\vspace{3mm}

\question[1] Which pair of lines is \textbf{parallel}?
\begin{choices}
    \choice $y = 2x + 1$ and $y = -2x + 1$
    \correctchoice $y = 3x - 5$ and $6x - 2y = 8$
    \choice $y = x + 4$ and $y = -x + 4$
    \choice $2x + 3y = 6$ and $3x + 2y = 6$
\end{choices}
\vspace{3mm}

\question[1] The point $(3, -4)$ lies on a circle centred at the origin. What is the radius?
\begin{choices}
    \choice $1$
    \choice $\sqrt{7}$
    \choice $4$
    \correctchoice $5$
\end{choices}
\vspace{3mm}

\question[1] What is the equation of the circle centred at the origin with radius $\sqrt{13}$?
\begin{choices}
    \choice $x^2 + y^2 = \sqrt{13}$
    \correctchoice $x^2 + y^2 = 13$
    \choice $(x-\sqrt{13})^2 + y^2 = 0$
    \choice $x^2 + y^2 = 169$
\end{choices}
\vspace{3mm}

\question[1] Triangle $ABC$ has vertices $A(0,0)$, $B(6,0)$, $C(6,8)$. What type of triangle is it?
\begin{choices}
    \choice Equilateral
    \choice Isosceles and acute
    \correctchoice Right triangle
    \choice Scalene acute
\end{choices}
\vspace{3mm}

\question[1] The median from vertex $A(2, 4)$ to side $BC$ where $B(0,0)$ and $C(4, 2)$ goes to midpoint:
\begin{choices}
    \choice $(2, 1)$
    \correctchoice $(2, 1)$
    \choice $(1, 2)$
    \choice $(3, 2)$
\end{choices}
\vspace{3mm}

\question[1] An altitude of a triangle is:
\begin{choices}
    \choice A segment from a vertex to the midpoint of the opposite side.
    \correctchoice A segment from a vertex perpendicular to the opposite side (or its extension).
    \choice The perpendicular bisector of a side.
    \choice A segment connecting the midpoints of two sides.
\end{choices}
\vspace{3mm}

\question[1] What is the slope of a line parallel to $4x - 2y + 7 = 0$?
\begin{choices}
    \choice $-2$
    \choice $4$
    \correctchoice $2$
    \choice $-4$
\end{choices}
\vspace{3mm}

\question[1] The perpendicular bisector of segment $PQ$ passes through:
\begin{choices}
    \choice Both $P$ and $Q$
    \choice Only the midpoint of $PQ$, not necessarily perpendicular
    \correctchoice The midpoint of $PQ$ at a right angle to $PQ$
    \choice One of the endpoints only
\end{choices}
\vspace{3mm}

\question[1] Which point lies on the circle $x^2 + y^2 = 25$?
\begin{choices}
    \choice $(3, 5)$
    \correctchoice $(4, -3)$
    \choice $(-5, 5)$
    \choice $(2, \sqrt{3})$
\end{choices}
\vspace{3mm}

\question[1] The equation of a line through $(1, 5)$ with slope $-2$ in standard form is:
\begin{choices}
    \choice $2x + y = 7$
    \choice $y = -2x + 7$
    \correctchoice $2x + y - 7 = 0$
    \choice $-2x + y + 3 = 0$
\end{choices}
\vspace{5mm}

%%%%%%%%%%%%%%%%%%%%%%%%%%%%%%%%%%%%%%%%%%%%%%%%%%%
\section*{Part B --- Short Answer (33 marks)}
%%%%%%%%%%%%%%%%%%%%%%%%%%%%%%%%%%%%%%%%%%%%%%%%%%%

\question \textbf{Triangle Side Lengths. (6 marks)}\\
Triangle $PQR$ has vertices $P(1, 2)$, $Q(7, 2)$, and $R(4, 6)$.

\begin{parts}
\part[3] Find the length of each side $PQ$, $QR$, and $PR$. Leave answers in exact (surd) form.

\begin{solution}[7cm]
\[
PQ = \sqrt{(7-1)^2 + (2-2)^2} = \sqrt{36} = 6
\]
\[
QR = \sqrt{(4-7)^2 + (6-2)^2} = \sqrt{9+16} = \sqrt{25} = 5
\]
\[
PR = \sqrt{(4-1)^2 + (6-2)^2} = \sqrt{9+16} = \sqrt{25} = 5
\]
\end{solution}

\part[2] Classify the triangle as equilateral, isosceles, or scalene. Justify your answer.

\begin{solution}[3cm]
Since $QR = PR = 5$ but $PQ = 6$, the triangle is \textbf{isosceles}.
\end{solution}

\part[1] Is the triangle right-angled? Justify.

\begin{solution}[3cm]
Check with Pythagorean theorem (longest side is $PQ = 6$):
$5^2 + 5^2 = 50 \neq 36 = 6^2$. \textbf{Not a right triangle.}
\end{solution}
\end{parts}

%---
\question \textbf{Triangle Line Segments. (6 marks)}\\
Triangle $ABC$ has vertices $A(0, 6)$, $B(-4, 0)$, and $C(4, 0)$.

\begin{parts}
\part[2] Find the equation of the \textbf{median} from $A$ to side $BC$.

\begin{solution}[5cm]
Midpoint of $BC$: $M = \left(\dfrac{-4+4}{2},\ \dfrac{0+0}{2}\right) = (0, 0)$.\\
Slope of $AM$: $m = \dfrac{0 - 6}{0 - 0}$ --- this is undefined (vertical line $x = 0$).\\
Equation of median: $\boxed{x = 0}$ (the $y$-axis).
\end{solution}

\part[2] Find the equation of the \textbf{altitude} from $B$ to side $AC$.

\begin{solution}[5cm]
Slope of $AC$: $m_{AC} = \dfrac{0 - 6}{4 - 0} = \dfrac{-6}{4} = -\dfrac{3}{2}$.\\
Slope of altitude (perpendicular to $AC$): $m = \dfrac{2}{3}$.\\
Altitude passes through $B(-4, 0)$:
\[
y - 0 = \frac{2}{3}(x + 4) \implies y = \frac{2}{3}x + \frac{8}{3}
\]
Or in standard form: $2x - 3y + 8 = 0$.
\end{solution}

\part[2] Find the equation of the \textbf{perpendicular bisector} of $AC$.

\begin{solution}[5cm]
Midpoint of $AC$: $N = \left(\dfrac{0+4}{2},\ \dfrac{6+0}{2}\right) = (2, 3)$.\\
Slope of $AC = -\dfrac{3}{2}$, so slope of perpendicular bisector $= \dfrac{2}{3}$.\\
Through $(2, 3)$:
\[
y - 3 = \frac{2}{3}(x - 2) \implies y = \frac{2}{3}x + \frac{5}{3}
\]
Or: $2x - 3y + 5 = 0$.
\end{solution}
\end{parts}

%---
\question \textbf{Parallel, Perpendicular, or Neither. (5 marks)}

\begin{parts}
\part[3] Determine whether each pair of lines is parallel, perpendicular, or neither.
\begin{enumerate}[(i)]
    \item $y = \dfrac{3}{4}x - 1$ \quad and \quad $4x - 3y + 6 = 0$
    \item $2x + 5y = 10$ \quad and \quad $5x - 2y = 4$
    \item $y = -3x + 2$ \quad and \quad $y = -3x - 7$
\end{enumerate}

\begin{solution}[6cm]
\begin{enumerate}[(i)]
    \item Rewrite $4x - 3y + 6 = 0 \Rightarrow y = \tfrac{4}{3}x + 2$, slope $= \tfrac{4}{3}$.\\
    First line slope $= \tfrac{3}{4}$. Product $= \tfrac{3}{4} \cdot \tfrac{4}{3} = 1 \neq -1$; not equal. \textbf{Neither.}
    \item Slope of $2x + 5y = 10$: $m_1 = -\tfrac{2}{5}$.\\
    Slope of $5x - 2y = 4$: $m_2 = \tfrac{5}{2}$.\\
    $m_1 \cdot m_2 = -\tfrac{2}{5} \cdot \tfrac{5}{2} = -1$. \textbf{Perpendicular.}
    \item Both have slope $-3$. \textbf{Parallel.}
\end{enumerate}
\end{solution}

\part[2] Find the equation of the line through $(3, -1)$ that is perpendicular to $y = \dfrac{3}{2}x + 4$.

\begin{solution}[4cm]
Slope of given line $= \dfrac{3}{2}$. Perpendicular slope $= -\dfrac{2}{3}$.\\
Through $(3, -1)$:
\[
y + 1 = -\frac{2}{3}(x - 3) \implies y = -\frac{2}{3}x + 1
\]
\end{solution}
\end{parts}

%---
\question \textbf{Circle Problems. (8 marks)}

\begin{parts}
\part[2] Determine whether each point lies on, inside, or outside the circle $x^2 + y^2 = 50$.
\begin{enumerate}[(i)]
    \item $A(5, 5)$
    \item $B(6, 4)$
\end{enumerate}

\begin{solution}[4cm]
\begin{enumerate}[(i)]
    \item $5^2 + 5^2 = 25 + 25 = 50$. Point $A$ lies \textbf{on} the circle.
    \item $6^2 + 4^2 = 36 + 16 = 52 > 50$. Point $B$ lies \textbf{outside} the circle.
\end{enumerate}
\end{solution}

\part[3] Find the points of intersection of the line $y = x + 2$ and the circle $x^2 + y^2 = 10$.

\begin{solution}[7cm]
Substitute $y = x + 2$ into $x^2 + y^2 = 10$:
\[
x^2 + (x+2)^2 = 10
\]
\[
x^2 + x^2 + 4x + 4 = 10
\]
\[
2x^2 + 4x - 6 = 0 \implies x^2 + 2x - 3 = 0
\]
\[
(x+3)(x-1) = 0 \implies x = -3 \text{ or } x = 1
\]
When $x = -3$: $y = -3 + 2 = -1$.\\
When $x = 1$: $y = 1 + 2 = 3$.\\
Intersection points: $\mathbf{(-3, -1)}$ and $\mathbf{(1, 3)}$.
\end{solution}

\part[3] A circle centred at the origin passes through $(-3, 4)$. Write its equation and determine whether $(5, 0)$ is on the circle.

\begin{solution}[5cm]
Radius $r = \sqrt{(-3)^2 + 4^2} = \sqrt{9 + 16} = 5$.\\
Equation: $x^2 + y^2 = 25$.\\
Check $(5, 0)$: $5^2 + 0^2 = 25$. Yes, $(5, 0)$ \textbf{lies on} the circle.
\end{solution}
\end{parts}

%---
\question \textbf{Coordinate Geometry Proof. (8 marks)}\\
Quadrilateral $ABCD$ has vertices $A(-3, 1)$, $B(1, 4)$, $C(5, 1)$, $D(1, -2)$.

\begin{parts}
\part[2] Find the midpoint of diagonal $AC$ and the midpoint of diagonal $BD$.

\begin{solution}[4cm]
Midpoint of $AC$: $\left(\dfrac{-3+5}{2},\ \dfrac{1+1}{2}\right) = (1, 1)$.\\
Midpoint of $BD$: $\left(\dfrac{1+1}{2},\ \dfrac{4+(-2)}{2}\right) = (1, 1)$.
\end{solution}

\part[2] What do your results from part (a) tell you about the diagonals?

\begin{solution}[3cm]
Both diagonals have the same midpoint $(1, 1)$, which means they \textbf{bisect each other}. Therefore $ABCD$ is a parallelogram.
\end{solution}

\part[2] Find the lengths of $AB$, $BC$, $CD$, and $DA$.

\begin{solution}[5cm]
$AB = \sqrt{(1+3)^2+(4-1)^2} = \sqrt{16+9} = 5$\\
$BC = \sqrt{(5-1)^2+(1-4)^2} = \sqrt{16+9} = 5$\\
$CD = \sqrt{(1-5)^2+(-2-1)^2} = \sqrt{16+9} = 5$\\
$DA = \sqrt{(-3-1)^2+(1+2)^2} = \sqrt{16+9} = 5$
\end{solution}

\part[2] Use your side lengths to prove that $ABCD$ is a rhombus.

\begin{solution}[3cm]
All four sides have equal length ($AB = BC = CD = DA = 5$). Since $ABCD$ is a parallelogram (diagonals bisect each other) with all sides equal, it is a \textbf{rhombus}.
\end{solution}
\end{parts}

\end{questions}
\end{document}
